\section{Limitations}\label{limitations}

Several limitations warrant discussion. The most important one, already flagged in Section~\ref{introduction}, deserves restating in full: our modeling framework is grounded in Earth-originating scenario typologies, drawing from narrative futures methodologies informed by contemporary geopolitical, ecological, and technological trends. These projections embed assumptions about governance structures, resource constraints, and recovery dynamics that may not extend to non-terrestrial civilizations or to those that do not follow Earth-like developmental pathways. While this grounding enables empirical plausibility, it limits extrapolation beyond our biosphere's trajectory \citep{Rubin2001-sk,Denning2011-fj}. The scenario ensemble is a structured sample of imagined Earth futures crafted via a ``Social Shaping of Technology'' approach to technology foresight which explicitly recognizes the ``co-shaping of technology and society'' \citep{Jorgensen2009-ss}, not a statistical sample of civilizations, so the specific duty-cycle values reported here should be read as illustrative of the mechanism rather than as estimates of a galactic distribution. Although the simulations are numerical, what we consider transferable is structural. Any civilization subject to resource depletion, exogenous hazards, and imperfect recovery, whatever the values of those parameters, will exhibit a duty cycle below unity, and the closed-form asymptotic expression (Supplementary Section~S6) makes explicit how that duty cycle depends on the underlying quantities. The claim we advance for extraterrestrial application is this structural one, not the particular numbers. Our simulations assume bounded planetary systems and do not account for galactic colonization, interstellar migration, or modes of expansion that could alter the spatial and temporal visibility of civilizations \citep{Walters1980-ev,Papagiannis1983-xv}. We also exclude the effects of cultural convergence, information saturation, or technological transcendence (e.g., the emergence of post-biological intelligence), all of which could suppress or qualitatively change technosignature output. Our operational definition of detectability (based on internal sociotechnical continuity) does not differentiate among specific technosignature classes (e.g., radio emissions, waste heat, megastructures), nor does it address detection thresholds or instrumental sensitivity. A full assessment of the Great Silence would additionally require the number and spatial distribution of civilizations, signal strengths, and survey durations, all of which are outside the scope of our model. The detectability conclusions of Section~\ref{technosignatures-implications} therefore concern one contributing mechanism among several, evaluated within an Earth-derived scenario set.

Additionally, our assumption that technological capacity increases linearly each year according to a fixed growth rate $r$ is a significant simplification. The many discontinuities implied by the scenarios (not only collapse points but also radical transformation points such as the emergence of artificial general intelligence or breakthroughs in energy harvesting) suggest that a linear growth model may not adequately capture the dynamics of technological change across all regimes. Two robustness results bound the consequences of this simplification. First, because $T(t)$ does not feed back on resource dynamics in the baseline model, all collapse-recovery metrics are strictly invariant to the functional form of growth (linear, exponential, or logistic), a property we verify numerically (Supplementary Section~S9); the growth law affects only the reported technology trajectories, not the duty cycles or collapse statistics. Second, when we introduce such feedback, coupling technology to the depletion rate (technology as efficiency) or to the recovery fraction (technology as resilience), moderate coupling strengths shift duty cycles by at most a few percent and largely preserve the scenario ranking (Supplementary Section~S9). The deeper limitation is therefore not the linearity of growth but the absence of strong technology-resource feedback in the baseline model. Moreover, the available resource stock $R_0$ is treated as an exogenous parameter, but in practice it is contingent on the technological capacity to identify and utilize resources. For example, a civilization that develops asteroid mining or nuclear fusion effectively increases its accessible resource base, fundamentally altering the growth-depletion dynamics. As discussed in Section~\ref{earth-implications}, \citet{Haqq-Misra2025-lum} reframe the Kardashev scale as a luminosity limit constrained by thermodynamic efficiency (exergy) and explore exploration-exploitation trade-offs whose distinction our linear growth model does not capture. Future work could incorporate exergy constraints and non-linear growth functions informed by their luminosity-limit model, as well as feedback between technological capacity and accessible resource stocks, to produce more physically grounded projections of civilizational trajectories. Such couplings could, for example, describe rebound dynamics such as the Jevons paradox \citep{Jevons1865-cq}, where efficiency gains can spur greater consumption, or Kuznets-type relationships between development and resource impact.
